\documentclass[a4paper]{article}
\usepackage[cm]{fullpage}
\usepackage{array}
\usepackage{longtable}
\usepackage[nodayofweek]{datetime}

\author{Mattia Colombo}
\title{Thesis Revisions \\\vspace{.5cm}
	\large Acoustic Information Retrieval for Interactive Sound Rendering in Virtual Environments}
\date{\today}


\newcolumntype{C}[1]{>{\centering}p{#1}}
\newcommand{\answer}[1]{\underline{\underline{#1}}}

\begin{document}
\maketitle

\begin{center}
	\begin{longtable}{|c|C{7.5cm}|C{7.5cm}|}
		\hline 
		\bf{No.} & \bf{Revision} & \bf{Action Taken} \tabularnewline
		\hline
		1. & Abstract revision: emphasise the research's originality and importance by clearly outlining the pipeline's purpose and concisely summarising the key findings for a stronger impact & Clarified ambiguity between AR and MR terms. Novel contributions explicitly stated in the abstract. Rewritten and refactored paragraphs to improve clarity and overall conciness of the abstract, fitting it into a single page. \tabularnewline
		\hline

		2. & Language Consistency: Ensure the dissertation uses British English consistently throughout. & Checked and corrected instances of grammar errors, spelling mistakes, and unclear sentences throughout the whole manuscript. \tabularnewline
		\hline
		
		3. & Terminology Standardisation: Maintain uniform terminology across the document, particularly when discussing AR, MR, and XR technologies. & Clarified the wrong use of the language and terminology around XR, MR, and AR across the relevant chapters and in the abstract. \tabularnewline
		\hline
		
		4. & Dataset Description: Provide a clear explanation of the dataset, including any limitations and the segmentation 
		methods applied. & Added methodology details on the dataset generation for the camera-based system to increase reproducibility and stated limitations of the trained models around their generalisability at the top of the Chapter. \tabularnewline
		\hline
		
		5. & Figure Citation and Organisation: Ensure that all figures are correctly cited and logically integrated within the text. Re-arrange figures where needed to improve the document’s flow. & All Figures are correctly referenced. Appropriate Figures were resized or eliminated to remove duplicates. \tabularnewline
		\hline

		6. & Avoid Figure Duplication: Ensure each figure is unique within the thesis and avoid repeating the same visuals across sections. & As per above, some Figures were eliminated to remove duplicates and correctly cross-referenced. Honouring the comments from the annotated PDF, repetitions were addressed throughout the thesis, adding cross-references where applicable. \tabularnewline
		\hline

		7. & Table 4.3 Clarity: Organise the results in Table 4.3 to enhance readability and prevent potential confusion. & The Table in question has been re-rendered, dividing it into two parts and making it clear that it shows two scenes. The caption has been updated to further clarify. All other Figures and Tables annotated in the PDF have been updated to include axes, labels, and better captions. \tabularnewline
		\hline

		8. & Figure 4.10 Explanation: Provide additional clarification on Figure 4.10, specifically regarding the predicted class labels for improved understanding. & The three Figures in questions have been grouped into one, using subfigures, adding a caption to clarify the process visualised by them. The figure is now standalone, and there are in-text references that explain the engineering process. \tabularnewline
		\hline
		
		9. & There are some writing issues in the thesis that need to be fixed. One of the examples are the sentences that and not well structured and do not make much sense (e.g. missing a verb) – see the comments in the PDF. & Improvements on the writing, paragraph structure, and sentences are made following all annotations provided on the PDF. This was carried out on all Chapters.\tabularnewline
		\hline

		10. & Section 1.1. Define and introduce presence, immersion, and embodiment, and then link them to audio, and then add the interaction and interfaces. & As per discussion, the term embodiment is not relevant to the work and has been removed from the manuscript. Sections 1.1.1 and 1.1.2 have been collapsed into one, making their writing more concise and improving their readability.\tabularnewline
		\hline 
		
		11. & Section 1.1.5 should be rewritten. There are some writing issues, but also it talks mostly about XR and briefly mentions the audio interaction system – this should be balanced better. & Following the guidance from the PDF and similarly to Revision 10, Sections 1.1.4 and 1.1.5 have been collapsed into one, making their writing more concise and improving their readability.\tabularnewline
		\hline 
		
		12. & Make the contributions (Section 1.7) a bit clearer, using more consistent terminology and potentially listing them chronologically. & As per discussions with the committee, Key Performance indicators were removed from Section 1.4. Section 1.7 has been improved, making the contributions concise and easier to read.\tabularnewline
		\hline 

		13. & Reconsider the sizes of the figures in the thesis. Some of them are unnecessarily too large. & This has been addressed as part of Revision 5.\tabularnewline
		\hline

		14. & I am slightly concerned about the repetition and redundancy in the thesis. Some concepts/definitions/explanations are repeated multiple times in the thesis which could have been avoided. It is good to link sections and chapters, but you do not need to explain same things multiple times. See the PDF for specific examples. & This has been addressed as part of Revisions 6 and 9. All specific examples annotated in the PDF are addressed.\tabularnewline
		\hline

		15. & Some details on the method and the setup are provided in the results (e.g. Section 4.2.4), although these should be mentioned earlier. & Section 4.2.4 and 4.3.3 have been edited, moving all methodology aspects to earlier Sections or any discussion about results to appropriate Discussion Sections. The whole Chapter has been improved to clarify unclear parts, as per recommendations provided.\tabularnewline
		\hline
		
		16. & I am slightly bothered with calling Section 5.2.8 Subjective Evaluation. A perceptual model CDPAM, trained on human judgements is used, but it is still algorithmic. This distinction is important because it shifts the evaluation from a truly subjective approach (direct human feedback) to an indirect or model-based approach that emulates human perceptual evaluation. A more precise term might be "perceptual model-based evaluation" or "model-based perceptual similarity assessment". This terminology would clarify that, while the model is trained on subjective data, the evaluation is not directly subjective but rather a proxy of subjective assessment. & The relevant questions in Chapter 5 have been updated to clarify that the data is extracted using Learned Similarity Metrics. Honouring the PDF annotations, the rest of the Chapter has been improved.\tabularnewline
		\hline

	\end{longtable}
\end{center}
\end{document}